\sectionheading{5.46 ব্যবহারিক Python Program অনুশীলন}

Programming শেখার সবচেয়ে কার্যকর উপায় হলো Problem বুঝে নিজে
Logic তৈরি করে Program লেখা।

\begin{center}

\textbf{Input নির্ধারণ}
\(\longrightarrow\)
\textbf{Processing/Condition}
\(\longrightarrow\)
\textbf{Code লেখা}
\(\longrightarrow\)
\textbf{Output যাচাই}

\end{center}

%================================================
% Strategy
%================================================

\booksubheading{Program লেখার সাধারণ কৌশল}

\begin{tabularx}
{\linewidth}
{
|>{\bfseries\RaggedRight\arraybackslash}p{0.48\linewidth}
|Y|
}

\hline
\rowcolor{softBlue}

\textbf{প্রশ্ন}
&
\textbf{কী নির্ধারণ করবে?}
\\
\hline

কী কী Data নিতে হবে?
&
Input
\\
\hline

কী Calculation বা Condition লাগবে?
&
Processing
\\
\hline

কী Result দেখাতে হবে?
&
Output
\\
\hline

\end{tabularx}

%================================================
% 1 Sum
%================================================

\booksubheading{১. দুটি সংখ্যার যোগফল}

\begin{codeoutputbox}{দুটি সংখ্যার Sum}

\noindent
\begin{minipage}[t]{0.57\linewidth}
\codeheading

\begin{english}
\ttfamily
a = int(input("Enter first number: "))\\
b = int(input("Enter second number: "))\\[0.5mm]
sum = a + b\\
print("Sum =", sum)
\end{english}

\end{minipage}
\hfill
\begin{minipage}[t]{0.39\linewidth}
\outputheading

\textbf{Input:}
\textenglish{\ttfamily 10, 20}

\medskip

\textbf{Output:}

\textenglish{\ttfamily Sum = 30}

\medskip

\textbf{IPO:}

Input \(\rightarrow\) \textenglish{\ttfamily a, b}

Processing \(\rightarrow\)
\textenglish{\ttfamily a + b}

Output \(\rightarrow\) Sum

\end{minipage}

\end{codeoutputbox}

%================================================
% 2 Rectangle
%================================================

\booksubheading{২. আয়তক্ষেত্রের ক্ষেত্রফল}

সূত্র—
\[
\text{Area} = \text{Length} \times \text{Width}
\]

\begin{codeoutputbox}{Rectangle Area}

\noindent
\begin{minipage}[t]{0.58\linewidth}
\codeheading

\begin{english}
\ttfamily
length = float(input("Enter length: "))\\
width = float(input("Enter width: "))\\[0.5mm]
area = length * width\\
print("Area =", area)
\end{english}

\end{minipage}
\hfill
\begin{minipage}[t]{0.38\linewidth}
\outputheading

\textbf{Input:}
\textenglish{\ttfamily 5, 4}

\medskip

\textbf{Output:}

\textenglish{\ttfamily Area = 20.0}

\medskip

\textenglish{\ttfamily float()}
ব্যবহার করায় Decimal Input-ও গ্রহণ করা যাবে।

\end{minipage}

\end{codeoutputbox}

%================================================
% 3 Average
%================================================

\booksubheading{৩. তিনটি সংখ্যার গড়}

\[
\text{Average}=\frac{a+b+c}{3}
\]

\begin{codeoutputbox}{Average নির্ণয়}

\noindent
\begin{minipage}[t]{0.59\linewidth}
\codeheading

\begin{english}
\ttfamily
a = float(input("Enter first number: "))\\
b = float(input("Enter second number: "))\\
c = float(input("Enter third number: "))\\[0.5mm]
avg = (a + b + c) / 3\\
print("Average =", avg)
\end{english}

\end{minipage}
\hfill
\begin{minipage}[t]{0.37\linewidth}
\outputheading

\textbf{Input:}

\textenglish{\ttfamily 10, 20, 30}

\medskip

\textbf{Output:}

\textenglish{\ttfamily Average = 20.0}

\medskip

তিনটি সংখ্যার যোগফলকে 3 দিয়ে ভাগ করা হয়েছে।

\end{minipage}

\end{codeoutputbox}

%================================================
% 4 Celsius Fahrenheit
%================================================

\booksubheading{৪. Celsius থেকে Fahrenheit}

সূত্র—

\[
F=\frac{9C}{5}+32
\]

\begin{codeoutputbox}{Temperature Conversion}

\noindent
\begin{minipage}[t]{0.56\linewidth}
\codeheading

\begin{english}
\ttfamily
c = float(input("Enter Celsius: "))\\[0.5mm]
f = (9 * c) / 5 + 32\\
print("Fahrenheit =", f)
\end{english}

\end{minipage}
\hfill
\begin{minipage}[t]{0.40\linewidth}
\outputheading

\textbf{Input:}
\textenglish{\ttfamily 25}

\medskip

\textbf{Output:}

\textenglish{\ttfamily Fahrenheit = 77.0}

\medskip

Formula অনুযায়ী Celsius-কে Fahrenheit-এ রূপান্তর করা হয়েছে।

\end{minipage}

\end{codeoutputbox}

%================================================
% 5 Even Odd
%================================================

\booksubheading{৫. সংখ্যা জোড় না বিজোড়}

\begin{codeoutputbox}{Even / Odd}

\noindent
\begin{minipage}[t]{0.54\linewidth}
\codeheading

\begin{english}
\ttfamily
n = int(input("Enter a number: "))\\[0.5mm]
if n \% 2 == 0:\\
\hspace*{4mm}print("Even")\\
else:\\
\hspace*{4mm}print("Odd")
\end{english}

\end{minipage}
\hfill
\begin{minipage}[t]{0.42\linewidth}
\outputheading

\textbf{Input:}
\textenglish{\ttfamily 8}

\medskip

\textbf{Output:}
\textenglish{\ttfamily Even}

\medskip

\textenglish{\ttfamily n \% 2 == 0}
হলে 2 দিয়ে ভাগশেষ 0।

তাই সংখ্যাটি Even।

\end{minipage}

\end{codeoutputbox}

\begin{rememberbox}

\begin{center}

\textenglish{\ttfamily n \% 2 == 0}
\(\longrightarrow\)
\textbf{Even}

\qquad

অন্যথায়
\(\longrightarrow\)
\textbf{Odd}

\end{center}

\end{rememberbox}

%================================================
% 6 Positive Negative
%================================================

\booksubheading{৬. সংখ্যা ধনাত্মক, ঋণাত্মক নাকি শূন্য}

\begin{codeoutputbox}{Positive / Negative / Zero}

\noindent
\begin{minipage}[t]{0.54\linewidth}
\codeheading

\begin{english}
\ttfamily
n = float(input("Enter a number: "))\\[0.5mm]
if n > 0:\\
\hspace*{4mm}print("Positive")\\
elif n < 0:\\
\hspace*{4mm}print("Negative")\\
else:\\
\hspace*{4mm}print("Zero")
\end{english}

\end{minipage}
\hfill
\begin{minipage}[t]{0.42\linewidth}
\outputheading

তিনটি সম্ভাবনা—

\medskip

\textenglish{\ttfamily n > 0}
\(\rightarrow\)
Positive

\medskip

\textenglish{\ttfamily n < 0}
\(\rightarrow\)
Negative

\medskip

অন্যথায়
\(\rightarrow\)
Zero

\medskip

তাই
\textenglish{\ttfamily if-elif-else}
ব্যবহার করা হয়েছে।

\end{minipage}

\end{codeoutputbox}

%================================================
% 7 Largest of two
%================================================

\booksubheading{৭. দুটি সংখ্যার মধ্যে বড় সংখ্যা}

\begin{codeoutputbox}{দুটি সংখ্যার তুলনা}

\noindent
\begin{minipage}[t]{0.57\linewidth}
\codeheading

\begin{english}
\ttfamily
a = int(input("Enter first number: "))\\
b = int(input("Enter second number: "))\\[0.5mm]
if a > b:\\
\hspace*{4mm}print(a, "is larger")\\
elif b > a:\\
\hspace*{4mm}print(b, "is larger")\\
else:\\
\hspace*{4mm}print("Both are equal")
\end{english}

\end{minipage}
\hfill
\begin{minipage}[t]{0.39\linewidth}
\outputheading

দুটি Number-এর মধ্যে—

\medskip

\textenglish{\ttfamily a > b}
হলে
\textenglish{\ttfamily a}
বড়।

\medskip

\textenglish{\ttfamily b > a}
হলে
\textenglish{\ttfamily b}
বড়।

\medskip

সমান হলে
\textenglish{\ttfamily else}
Execute হয়।

\end{minipage}

\end{codeoutputbox}

%================================================
% 8 Largest three
%================================================

\booksubheading{৮. তিনটি সংখ্যার মধ্যে বৃহত্তম সংখ্যা}

\begin{codeoutputbox}{Largest of Three}

\noindent
\begin{minipage}[t]{0.59\linewidth}
\codeheading

\begin{english}
\ttfamily
a = int(input("Enter first number: "))\\
b = int(input("Enter second number: "))\\
c = int(input("Enter third number: "))\\[0.5mm]
if a >= b and a >= c:\\
\hspace*{4mm}largest = a\\
elif b >= a and b >= c:\\
\hspace*{4mm}largest = b\\
else:\\
\hspace*{4mm}largest = c\\[0.5mm]
print("Largest =", largest)
\end{english}

\end{minipage}
\hfill
\begin{minipage}[t]{0.37\linewidth}
\outputheading

\textbf{Input:}

\textenglish{\ttfamily 10, 25, 15}

\medskip

\textbf{Output:}

\textenglish{\ttfamily Largest = 25}

\medskip

প্রতিটি Number-কে বাকি দুইটির সঙ্গে তুলনা করা হয়েছে।

\end{minipage}

\end{codeoutputbox}

%================================================
% 9 Grade
%================================================

\booksubheading{৯. Marks থেকে Grade নির্ণয়}

\begin{codeoutputbox}{Grade নির্ণয়}

\noindent
\begin{minipage}[t]{0.53\linewidth}
\codeheading

\begin{english}
\ttfamily
marks = int(input("Enter marks: "))\\[0.5mm]
if marks >= 80:\\
\hspace*{4mm}print("A+")\\
elif marks >= 70:\\
\hspace*{4mm}print("A")\\
elif marks >= 60:\\
\hspace*{4mm}print("A-")\\
elif marks >= 50:\\
\hspace*{4mm}print("B")\\
elif marks >= 40:\\
\hspace*{4mm}print("C")\\
elif marks >= 33:\\
\hspace*{4mm}print("D")\\
else:\\
\hspace*{4mm}print("F")
\end{english}

\end{minipage}
\hfill
\begin{minipage}[t]{0.43\linewidth}
\outputheading

Conditionগুলো উচ্চ Marks থেকে নিচের দিকে সাজানো হয়েছে।

\medskip

যেমন
\textenglish{\ttfamily marks = 75}
হলে—

\begin{center}
\textenglish{\ttfamily A}
\end{center}

Print হবে।

\medskip

একাধিক Range পরীক্ষা করতে
\textenglish{\ttfamily if-elif-else}
ব্যবহার করা হয়েছে।

\end{minipage}

\end{codeoutputbox}

%================================================
% 10 1 to n
%================================================

\booksubheading{১০. ১ থেকে n পর্যন্ত সংখ্যা Print}

\begin{codeoutputbox}{১ থেকে n}

\noindent
\begin{minipage}[t]{0.52\linewidth}
\codeheading

\begin{english}
\ttfamily
n = int(input("Enter n: "))\\[0.5mm]
for i in range(1, n + 1):\\
\hspace*{4mm}print(i)
\end{english}

\end{minipage}
\hfill
\begin{minipage}[t]{0.44\linewidth}
\outputheading

যদি
\textenglish{\ttfamily n = 5},
Output—

\begin{english}
\ttfamily
1\\
2\\
3\\
4\\
5
\end{english}

\textenglish{\ttfamily range()}
-এর Stop Value Excluded, তাই
\textenglish{\ttfamily n + 1}।

\end{minipage}

\end{codeoutputbox}

%================================================
% 11 Sum 1 to n
%================================================

\booksubheading{১১. ১ থেকে n পর্যন্ত সংখ্যার যোগফল}

\begin{codeoutputbox}{Accumulator ব্যবহার}

\noindent
\begin{minipage}[t]{0.53\linewidth}
\codeheading

\begin{english}
\ttfamily
n = int(input("Enter n: "))\\
total = 0\\[0.5mm]
for i in range(1, n + 1):\\
\hspace*{4mm}total += i\\[0.5mm]
print("Sum =", total)
\end{english}

\end{minipage}
\hfill
\begin{minipage}[t]{0.43\linewidth}
\outputheading

\textbf{Input:}
\textenglish{\ttfamily 5}

\medskip

\textbf{Output:}

\textenglish{\ttfamily Sum = 15}

\medskip

\textenglish{\ttfamily total}
প্রতিটি Iteration-এর Sum জমা রাখছে।

তাই এটি \textbf{Accumulator}।

\end{minipage}

\end{codeoutputbox}

\begin{tabularx}
{\linewidth}
{
|>{\bfseries\centering\arraybackslash}p{0.30\linewidth}
|>{\centering\arraybackslash}Y|
}

\hline
\rowcolor{softBlue}
\textbf{\textenglish{\ttfamily i}}
&
\textbf{\textenglish{\ttfamily total}}
\\
\hline
1 & 1 \\
\hline
2 & 3 \\
\hline
3 & 6 \\
\hline
4 & 10 \\
\hline
5 & 15 \\
\hline

\end{tabularx}

%================================================
% 12 Multiplication table
%================================================

\booksubheading{১২. গুণের নামতা}

\begin{codeoutputbox}{Multiplication Table}

\noindent
\begin{minipage}[t]{0.55\linewidth}
\codeheading

\begin{english}
\ttfamily
n = int(input("Enter a number: "))\\[0.5mm]
for i in range(1, 11):\\
\hspace*{4mm}print(n, "x", i, "=", n * i)
\end{english}

\end{minipage}
\hfill
\begin{minipage}[t]{0.41\linewidth}
\outputheading

যদি
\textenglish{\ttfamily n = 5}—

\begin{english}
\ttfamily
5 x 1 = 5\\
5 x 2 = 10\\
5 x 3 = 15\\
...\\
5 x 10 = 50
\end{english}

Loop 1 থেকে 10 পর্যন্ত চলে।

\end{minipage}

\end{codeoutputbox}

%================================================
% 13 Factorial
%================================================

\booksubheading{১৩. Factorial নির্ণয়}

\[
n! = 1 \times 2 \times 3 \times \cdots \times n
\]

\begin{codeoutputbox}{Factorial}

\noindent
\begin{minipage}[t]{0.53\linewidth}
\codeheading

\begin{english}
\ttfamily
n = int(input("Enter a number: "))\\
fact = 1\\[0.5mm]
for i in range(1, n + 1):\\
\hspace*{4mm}fact *= i\\[0.5mm]
print("Factorial =", fact)
\end{english}

\end{minipage}
\hfill
\begin{minipage}[t]{0.43\linewidth}
\outputheading

\textbf{Input:}
\textenglish{\ttfamily 5}

\medskip

\textbf{Output:}

\textenglish{\ttfamily Factorial = 120}

\medskip

Multiplication জমা করার জন্য শুরুতে—

\begin{center}
\textenglish{\ttfamily fact = 1}
\end{center}

নেওয়া হয়েছে।

\end{minipage}

\end{codeoutputbox}

\begin{rememberbox}

Factorial-এ
\textenglish{\ttfamily fact = 0}
নিলে প্রতিটি Multiplication-এর Result 0 হয়ে যাবে।

তাই Initial Value সাধারণত
\textenglish{\ttfamily 1}।

\end{rememberbox}

%================================================
% 14 Prime
%================================================

\booksubheading{১৪. একটি সংখ্যা Prime কি না}

Prime Number-এর ঠিক দুটি Positive Divisor থাকে—
\textenglish{\ttfamily 1}
এবং সংখ্যাটি নিজে।

\begin{codeoutputbox}{Prime Number Test}

\noindent
\begin{minipage}[t]{0.55\linewidth}
\codeheading

\begin{english}
\ttfamily
n = int(input("Enter a number: "))\\
is\_prime = True\\[0.5mm]
if n < 2:\\
\hspace*{4mm}is\_prime = False\\
else:\\
\hspace*{4mm}for i in range(2, n):\\
\hspace*{8mm}if n \% i == 0:\\
\hspace*{12mm}is\_prime = False\\
\hspace*{12mm}break\\[0.5mm]
if is\_prime:\\
\hspace*{4mm}print("Prime")\\
else:\\
\hspace*{4mm}print("Not Prime")
\end{english}

\end{minipage}
\hfill
\begin{minipage}[t]{0.41\linewidth}
\outputheading

\textbf{Input:}
\textenglish{\ttfamily 7}

\medskip

\textbf{Output:}
\textenglish{\ttfamily Prime}

\medskip

2 থেকে
\textenglish{\ttfamily n-1}
পর্যন্ত কোনো Number দিয়ে নিঃশেষে ভাগ যায় কি না দেখা হয়।

\medskip

Divisor পাওয়া গেলে
\textenglish{\ttfamily break}
দিয়ে Loop বন্ধ হয়।

\end{minipage}

\end{codeoutputbox}

%================================================
% 15 Fibonacci
%================================================

\booksubheading{১৫. Fibonacci Series}

Fibonacci Series-এ প্রতিটি নতুন সংখ্যা আগের দুই সংখ্যার যোগফল।

\begin{center}
\textenglish{\ttfamily 0, 1, 1, 2, 3, 5, 8, ...}
\end{center}

\begin{codeoutputbox}{Fibonacci Series}

\noindent
\begin{minipage}[t]{0.53\linewidth}
\codeheading

\begin{english}
\ttfamily
n = int(input("How many terms: "))\\
a = 0\\
b = 1\\[0.5mm]
for i in range(n):\\
\hspace*{4mm}print(a, end=" ")\\
\hspace*{4mm}a, b = b, a + b
\end{english}

\end{minipage}
\hfill
\begin{minipage}[t]{0.43\linewidth}
\outputheading

যদি
\textenglish{\ttfamily n = 7}—

\begin{center}
\textenglish{\ttfamily 0 1 1 2 3 5 8}
\end{center}

প্রতি Iteration-এ—

\begin{center}
\textenglish{\ttfamily a, b = b, a + b}
\end{center}

দিয়ে পরবর্তী দুইটি Value Update হয়।

\end{minipage}

\end{codeoutputbox}

%================================================
% 16 Digit sum
%================================================

\booksubheading{১৬. একটি সংখ্যার Digit-এর যোগফল}

উদাহরণ—

\[
1234 \longrightarrow 1+2+3+4=10
\]

\begin{codeoutputbox}{Sum of Digits}

\noindent
\begin{minipage}[t]{0.53\linewidth}
\codeheading

\begin{english}
\ttfamily
n = int(input("Enter a number: "))\\
total = 0\\[0.5mm]
while n > 0:\\
\hspace*{4mm}digit = n \% 10\\
\hspace*{4mm}total += digit\\
\hspace*{4mm}n //= 10\\[0.5mm]
print("Sum of digits =", total)
\end{english}

\end{minipage}
\hfill
\begin{minipage}[t]{0.43\linewidth}
\outputheading

\textbf{Input:}
\textenglish{\ttfamily 1234}

\medskip

\textbf{Output:}

\textenglish{\ttfamily Sum of digits = 10}

\medskip

\textenglish{\ttfamily n \% 10}
\(\rightarrow\)
শেষ Digit

\medskip

\textenglish{\ttfamily n // 10}
\(\rightarrow\)
শেষ Digit বাদ

\end{minipage}

\end{codeoutputbox}

%================================================
% 17 Reverse
%================================================

\booksubheading{১৭. সংখ্যা Reverse করা}

\begin{codeoutputbox}{Number Reverse}

\noindent
\begin{minipage}[t]{0.53\linewidth}
\codeheading

\begin{english}
\ttfamily
n = int(input("Enter a number: "))\\
reverse = 0\\[0.5mm]
while n > 0:\\
\hspace*{4mm}digit = n \% 10\\
\hspace*{4mm}reverse = reverse * 10 + digit\\
\hspace*{4mm}n //= 10\\[0.5mm]
print("Reverse =", reverse)
\end{english}

\end{minipage}
\hfill
\begin{minipage}[t]{0.43\linewidth}
\outputheading

\textbf{Input:}
\textenglish{\ttfamily 1234}

\medskip

\textbf{Output:}

\textenglish{\ttfamily Reverse = 4321}

\medskip

প্রতিবার শেষ Digit নিয়ে
\textenglish{\ttfamily reverse}
Variable-এর ডান পাশে যোগ করা হয়েছে।

\end{minipage}

\end{codeoutputbox}

%================================================
% 18 Palindrome
%================================================

\booksubheading{১৮. Palindrome Number পরীক্ষা}

যে Number উল্টো করলেও একই থাকে তাকে
\textbf{Palindrome Number} বলে।

যেমন—
\textenglish{\ttfamily 121, 1331}।

\begin{codeoutputbox}{Palindrome Test}

\noindent
\begin{minipage}[t]{0.55\linewidth}
\codeheading

\begin{english}
\ttfamily
n = int(input("Enter a number: "))\\
original = n\\
reverse = 0\\[0.5mm]
while n > 0:\\
\hspace*{4mm}digit = n \% 10\\
\hspace*{4mm}reverse = reverse * 10 + digit\\
\hspace*{4mm}n //= 10\\[0.5mm]
if original == reverse:\\
\hspace*{4mm}print("Palindrome")\\
else:\\
\hspace*{4mm}print("Not Palindrome")
\end{english}

\end{minipage}
\hfill
\begin{minipage}[t]{0.41\linewidth}
\outputheading

\textbf{Input:}
\textenglish{\ttfamily 121}

\medskip

\textbf{Output:}

\textenglish{\ttfamily Palindrome}

\medskip

প্রথমে Original Number রাখা হয়েছে।

তারপর Reverse বের করে—

\begin{center}
\textenglish{\ttfamily original == reverse}
\end{center}

পরীক্ষা করা হয়েছে।

\end{minipage}

\end{codeoutputbox}

%================================================
% 19 Pattern
%================================================

\booksubheading{১৯. Increasing Star Pattern}

\begin{codeoutputbox}{Star Pattern}

\noindent
\begin{minipage}[t]{0.53\linewidth}
\codeheading

\begin{english}
\ttfamily
for i in range(1, 5):\\
\hspace*{4mm}for j in range(i):\\
\hspace*{8mm}print("*", end=" ")\\
\hspace*{4mm}print()
\end{english}

\end{minipage}
\hfill
\begin{minipage}[t]{0.43\linewidth}
\outputheading

\begin{english}
\ttfamily
*\\
* *\\
* * *\\
* * * *
\end{english}

\medskip

Outer Loop
\(\rightarrow\)
Row

\medskip

Inner Loop
\(\rightarrow\)
প্রতি Row-এর Star

\end{minipage}

\end{codeoutputbox}

%================================================
% 20 List sum
%================================================

\booksubheading{২০. List-এর সব সংখ্যার যোগফল}

\begin{codeoutputbox}{List Sum}

\noindent
\begin{minipage}[t]{0.53\linewidth}
\codeheading

\begin{english}
\ttfamily
numbers = [10, 20, 30, 40]\\
total = 0\\[0.5mm]
for x in numbers:\\
\hspace*{4mm}total += x\\[0.5mm]
print("Sum =", total)
\end{english}

\end{minipage}
\hfill
\begin{minipage}[t]{0.43\linewidth}
\outputheading

\textbf{Output:}

\begin{center}
\textenglish{\ttfamily Sum = 100}
\end{center}

List-এর প্রতিটি Item
\textenglish{\ttfamily x}
-এর মধ্যে এসেছে এবং
\textenglish{\ttfamily total}
-এ যোগ হয়েছে।

\end{minipage}

\end{codeoutputbox}

%================================================
% 21 Largest in list
%================================================

\booksubheading{২১. List-এর বৃহত্তম সংখ্যা}

\begin{codeoutputbox}{Largest Item}

\noindent
\begin{minipage}[t]{0.55\linewidth}
\codeheading

\begin{english}
\ttfamily
numbers = [12, 45, 8, 30]\\
largest = numbers[0]\\[0.5mm]
for x in numbers:\\
\hspace*{4mm}if x > largest:\\
\hspace*{8mm}largest = x\\[0.5mm]
print("Largest =", largest)
\end{english}

\end{minipage}
\hfill
\begin{minipage}[t]{0.41\linewidth}
\outputheading

\textbf{Output:}

\begin{center}
\textenglish{\ttfamily Largest = 45}
\end{center}

প্রথম Item-কে শুরুতে Largest ধরা হয়েছে।

তারপর প্রতিটি Item-এর সঙ্গে তুলনা করে
\textenglish{\ttfamily largest}
Update হয়েছে।

\end{minipage}

\end{codeoutputbox}

%================================================
% 22 Even count
%================================================

\booksubheading{২২. List-এ Even Number গণনা}

\begin{codeoutputbox}{Counter ব্যবহার}

\noindent
\begin{minipage}[t]{0.55\linewidth}
\codeheading

\begin{english}
\ttfamily
numbers = [10, 15, 20, 27, 30]\\
count = 0\\[0.5mm]
for x in numbers:\\
\hspace*{4mm}if x \% 2 == 0:\\
\hspace*{8mm}count += 1\\[0.5mm]
print("Even count =", count)
\end{english}

\end{minipage}
\hfill
\begin{minipage}[t]{0.41\linewidth}
\outputheading

\textbf{Output:}

\begin{center}
\textenglish{\ttfamily Even count = 3}
\end{center}

Even Itemগুলো—

\begin{center}
\textenglish{\ttfamily 10, 20, 30}
\end{center}

প্রতিবার Even পাওয়া গেলে
\textenglish{\ttfamily count}
এক বাড়ে।

\end{minipage}

\end{codeoutputbox}

\begin{rememberbox}

\textbf{Accumulator}
\(\longrightarrow\)
Value যোগ/জমা করে।

\qquad

\textbf{Counter}
\(\longrightarrow\)
কতবার ঘটনা ঘটেছে তা গণনা করে।

\end{rememberbox}

%================================================
% 23 Dictionary
%================================================

\booksubheading{২৩. Dictionary থেকে Student-এর তথ্য দেখানো}

\begin{codeoutputbox}{Student Record}

\noindent
\begin{minipage}[t]{0.55\linewidth}
\codeheading

\begin{english}
\ttfamily
student = \{\\
\hspace*{4mm}"name": "Rafi",\\
\hspace*{4mm}"marks": 85\\
\}\\[0.5mm]
print("Name:", student["name"])\\
print("Marks:", student["marks"])
\end{english}

\end{minipage}
\hfill
\begin{minipage}[t]{0.41\linewidth}
\outputheading

\begin{english}
\ttfamily
Name: Rafi\\
Marks: 85
\end{english}

\medskip

Dictionary-তে Position-এর পরিবর্তে
\textbf{Key} ব্যবহার করে Value Access করা হয়েছে।

\end{minipage}

\end{codeoutputbox}

%================================================
% 24 Function addition
%================================================

\booksubheading{২৪. Function দিয়ে দুটি সংখ্যার যোগ}

\begin{codeoutputbox}{Function দিয়ে Sum}

\noindent
\begin{minipage}[t]{0.58\linewidth}
\codeheading

\begin{english}
\ttfamily
def add(a, b):\\
\hspace*{4mm}return a + b\\[0.5mm]
x = int(input("Enter first number: "))\\
y = int(input("Enter second number: "))\\[0.5mm]
result = add(x, y)\\
print("Sum =", result)
\end{english}

\end{minipage}
\hfill
\begin{minipage}[t]{0.38\linewidth}
\outputheading

\textbf{Input:}

\textenglish{\ttfamily 10, 20}

\medskip

\textbf{Output:}

\textenglish{\ttfamily Sum = 30}

\medskip

\textbf{Contract:}

Input
\(\rightarrow\)
\textenglish{\ttfamily a, b}

Processing
\(\rightarrow\)
\textenglish{\ttfamily a+b}

Output
\(\rightarrow\)
Sum

\end{minipage}

\end{codeoutputbox}

%================================================
% 25 Function even odd
%================================================

\booksubheading{২৫. Function দিয়ে Even/Odd পরীক্ষা}

\begin{codeoutputbox}{Even/Odd Function}

\noindent
\begin{minipage}[t]{0.57\linewidth}
\codeheading

\begin{english}
\ttfamily
def check\_even\_odd(n):\\
\hspace*{4mm}if n \% 2 == 0:\\
\hspace*{8mm}return "Even"\\
\hspace*{4mm}else:\\
\hspace*{8mm}return "Odd"\\[0.5mm]
number = int(input("Enter a number: "))\\
print(check\_even\_odd(number))
\end{english}

\end{minipage}
\hfill
\begin{minipage}[t]{0.39\linewidth}
\outputheading

\textbf{Input:}
\textenglish{\ttfamily 7}

\medskip

\textbf{Output:}

\textenglish{\ttfamily Odd}

\medskip

Function Number গ্রহণ করে Condition পরীক্ষা করেছে এবং
\textenglish{\ttfamily "Even"}
বা
\textenglish{\ttfamily "Odd"}
Return করেছে।

\end{minipage}

\end{codeoutputbox}

%================================================
% Multiple solutions
%================================================

\booksubheading{একই Problem একাধিকভাবে সমাধান}

Programming-এ একই Problem-এর একাধিক সঠিক Solution থাকতে পারে।

উদাহরণস্বরূপ তিনটি সংখ্যার বৃহত্তম Value Built-in Function দিয়েও—

\begin{codeoutputbox}{Built-in \textenglish{\ttfamily max()} ব্যবহার}

\noindent
\begin{minipage}[t]{0.50\linewidth}
\codeheading

\begin{english}
\ttfamily
a = 10\\
b = 25\\
c = 15\\[0.5mm]
largest = max(a, b, c)\\
print(largest)
\end{english}

\end{minipage}
\hfill
\begin{minipage}[t]{0.46\linewidth}
\outputheading

\begin{center}
\textenglish{\ttfamily 25}
\end{center}

একই Problem আগে
\textenglish{\ttfamily if-elif-else}
দিয়েও সমাধান করা হয়েছে।

\medskip

একই Problem-এর একাধিক সঠিক Algorithm থাকতে পারে।

\end{minipage}

\end{codeoutputbox}

\begin{rememberbox}

পরীক্ষা ও শেখার ক্ষেত্রে শুধু Built-in Function-এর উপর নির্ভর না করে
\textbf{Condition, Loop, Data Structure এবং Function ব্যবহার করে Logic তৈরি করতে পারা}
গুরুত্বপূর্ণ।

\end{rememberbox}

%================================================
% Shortcut
%================================================

\booksubheading{Practical Program Solve করার Shortcut}

Problem দেখলে—

\begin{tabularx}
{\linewidth}
{
|>{\bfseries\centering\arraybackslash}p{0.13\linewidth}
|Y|
}

\hline
\rowcolor{softBlue}

\textbf{ধাপ}
&
\textbf{প্রশ্ন}
\\
\hline

১
&
Input কী?
\\
\hline

২
&
Formula বা Condition কী?
\\
\hline

৩
&
Loop লাগবে কি?
\\
\hline

৪
&
Counter বা Accumulator লাগবে কি?
\\
\hline

৫
&
Output কী হবে?
\\
\hline

\end{tabularx}

%================================================
% Problem mapping
%================================================

\booksubheading{কোন Problem-এ কোন ধারণা?}

\begin{tabularx}
{\linewidth}
{
|>{\bfseries\RaggedRight\arraybackslash}p{0.46\linewidth}
|Y|
}

\hline
\rowcolor{softBlue}

\textbf{Problem}
&
\textbf{প্রধান ধারণা}
\\
\hline

যোগ, গড়, ক্ষেত্রফল
&
Operator
\\
\hline

Even/Odd
&
\textenglish{\ttfamily \%} + Conditional
\\
\hline

Positive/Negative
&
\textenglish{\ttfamily if-elif-else}
\\
\hline

Largest
&
Conditional
\\
\hline

Grade
&
Multiple
\textenglish{\ttfamily elif}
\\
\hline

১ থেকে n
&
\textenglish{\ttfamily for}
Loop
\\
\hline

Sum
&
Accumulator
\\
\hline

নামতা
&
Loop
\\
\hline

Factorial
&
Loop + Multiplication
\\
\hline

Prime
&
Loop + Condition +
\textenglish{\ttfamily break}
\\
\hline

Fibonacci
&
Variable Update
\\
\hline

Digit Sum
&
\textenglish{\ttfamily while + \% + //}
\\
\hline

Reverse
&
\textenglish{\ttfamily while}
Loop
\\
\hline

Palindrome
&
Reverse + Condition
\\
\hline

Pattern
&
Nested Loop
\\
\hline

List Sum
&
List + Loop
\\
\hline

Largest in List
&
Traversal + Comparison
\\
\hline

Even Count
&
Counter
\\
\hline

Student Record
&
Dictionary
\\
\hline

Reusable Calculation
&
Function
\\
\hline

\end{tabularx}

%================================================
% Common Mistakes
%================================================

\booksubheading{সাধারণ ভুল}

\begin{rememberbox}

\begin{itemize}

\item
\textenglish{\ttfamily input()}
থেকে পাওয়া String-কে প্রয়োজনমতো
\textenglish{\ttfamily int()}
বা
\textenglish{\ttfamily float()}
-এ Convert না করা।

\item Assignment
\textenglish{\ttfamily =}
এবং Comparison
\textenglish{\ttfamily ==}
গুলিয়ে ফেলা।

\item Even/Odd পরীক্ষায়
\textenglish{\ttfamily \%}
-এর পরিবর্তে
\textenglish{\ttfamily /}
ব্যবহার করা।

\item
\textenglish{\ttfamily range()}
-এর Stop Value Excluded—এটি ভুলে যাওয়া।

\item Accumulator-এর Initial Value ভুল দেওয়া।

\item Factorial-এ
\textenglish{\ttfamily fact = 0}
দেওয়া।

\item
\textenglish{\ttfamily while}
Loop-এ Variable Update না করা।

\item Nested Loop-এ Indentation ভুল করা।

\item Prime Test-এ Divisor পাওয়া যাওয়ার পরও অপ্রয়োজনীয় Loop চালানো।

\item List-এর প্রথম Index
\textenglish{\ttfamily 0}
—এটি ভুলে যাওয়া।

\item Function Define করলেও Call না করা।

\item
\textenglish{\ttfamily print()}
ও
\textenglish{\ttfamily return}
-এর পার্থক্য ভুলে যাওয়া।

\end{itemize}

\end{rememberbox}

%================================================
% Final Reminder
%================================================

\begin{rememberbox}

\textbf{\color{deepOrange}মনে রাখবে}

\begin{tabularx}
{\linewidth}
{
>{\bfseries\RaggedRight\arraybackslash}p{0.25\linewidth}
Y
>{\bfseries\RaggedRight\arraybackslash}p{0.25\linewidth}
Y
}

Input
&
\textenglish{\ttfamily input()}
&
Integer Input
&
\textenglish{\ttfamily int(input())}
\\[0.5mm]

Decimal Input
&
\textenglish{\ttfamily float(input())}
&
Remainder
&
\textenglish{\ttfamily \%}
\\[0.5mm]

Integer Division
&
\textenglish{\ttfamily //}
&
Decision
&
\textenglish{\ttfamily if / elif / else}
\\[0.5mm]

নির্দিষ্টবার Loop
&
\textenglish{\ttfamily for}
&
Condition-based Loop
&
\textenglish{\ttfamily while}
\\[0.5mm]

Sum জমা
&
Accumulator
&
সংখ্যা গণনা
&
Counter
\\[0.5mm]

Loop বন্ধ
&
\textenglish{\ttfamily break}
&
Pattern
&
Nested Loop
\\[0.5mm]

একাধিক Data
&
List / Dictionary
&
Reusable Logic
&
Function
\\

\end{tabularx}

\medskip

\begin{center}

\textbf{Problem বুঝো}
\(\longrightarrow\)
\textbf{Input চিহ্নিত করো}
\(\longrightarrow\)
\textbf{Logic/Formula তৈরি করো}

\medskip

\(\longrightarrow\)
\textbf{Code লেখো}
\(\longrightarrow\)
\textbf{Dry Run করো}
\(\longrightarrow\)
\textbf{Output যাচাই করো}

\end{center}

\end{rememberbox}

\small

%================================================
% Main Title
%================================================
